% ==================== 附录 ====================
\section{附录：数据、脚本与操作记录}

\subsection{数据目录与脚本}

本次作业用到的文件都在 \texttt{X-Ray/作业二/} 下：

\begin{itemize}[leftmargin=2.2em]
  \item \texttt{资料/原数据/200.txt--900.txt}：仪器导出的原始 txt，GBK 编码，前 9 行为元数据；
  \item \texttt{资料/原数据/ready/}：\pyA 输出的原始强度，三行表头 + Tab 分隔；
  \item \texttt{资料/原数据/ready\_norm/}：同一脚本输出的归一化强度（每个样品除以自己的最强峰再乘 100）；
  \item \texttt{资料/处理后数据/200.txt--900.txt}：整理后的成品数据，内容与 \texttt{ready\_norm/} 一致，方便直接取用；
  \item \texttt{资料/图片/}：本次 Origin 和 Jade 操作的全部截图，正文图~\ref{fig:data}--\ref{fig:origin-step6}、图~\ref{fig:jade}--\ref{fig:pdf-step6} 都取自这里；
  \item \texttt{脚本/make\_origin\_ready.py}：把原始 txt 转成 Origin 可导入的三行表头格式；
  \item \texttt{脚本/plot\_xrd\_multi.py}：读取 \texttt{ready\_norm/}，绘制图~\ref{fig:xrd} 的多曲线 + PDF 卡片图。
\end{itemize}

两个脚本的运行方式如下（在作业二根目录下）：

\begin{center}
\texttt{python .\textbackslash 脚本\textbackslash make\_origin\_ready.py 资料\textbackslash 原数据 --norm}\\
\vspace{2pt}
\texttt{python .\textbackslash 脚本\textbackslash plot\_xrd\_multi.py}
\end{center}

\pyA 只依赖 Python 标准库；\pyB 需要 numpy 和 matplotlib。两个脚本都只读原始数据，输出写到 \texttt{ready/}、\texttt{ready\_norm/} 和 \texttt{课程论文/assets/xrd/}。

\subsection{操作记录的几点补充}

PDF 卡片层的位置没有唯一标准，第 4 节给的数值只是本次的起点，换页面尺寸时要重新调。判断标准只有一条：卡片层的 X 轴与主图同为 20$^\circ$--80$^\circ$，卡片不盖住实验曲线。

换一批数据时，只要仪器导出的 txt 格式不变，直接对目录运行 \pyA 即可，脚本不用改。9 行元数据里的波长、步宽、角度范围、日期、最强峰位置和样品编号会写进 Comments 行，与数据一起保留。
